% Source-derived chapter 03: The prime to the characteristic aspects of the main result
% Original source: purity-flat-cohomology
\chapter{The prime to the characteristic aspects of the main result}
\label{purity-flat-cohomology-chapter-03}
\source{purity-flat-cohomology}

\begin{remark}[Source section]
\label{purity-flat-cohomology-chapter-03-source}
\source{purity-flat-cohomology}
\source{purity-flat-cohomology}
This chapter reproduces the corresponding section of the retrieved
LaTeX source, including its definitions, statements, examples, and
proofs. It has not yet been translated into Lean.
\notready
\end{remark}

% The source section title was: The prime to the characteristic aspects of the main result
For arguing our purity results, the first task is to dispose of the cases when the order of the coefficients is invertible. For this, we first give a new, perfectoid-based proof of the Gabber--Thomason purity for \'{e}tale cohomology of regular rings in \S\ref{Gabber-Thomason}. We then use it in \S\ref{et-depth-geom-depth} to deduce purity for \'{e}tale cohomology in the general singular case via a local Lefschetz style theorem. In \S\ref{nonabelian}, we present a nonabelian analogue of the results of \S\ref{et-depth-geom-depth}: a generalization of the Zariski--Nagata purity theorem.


%We accomplish this in  %below %, we deduce the prime to the characteristic aspects of \Cref{ell-semipurity-ann} (and hence also those of \Cref{main}) from 
%by deducing them form the purity for \'{e}tale cohomology due to Gabber--Thomason. We  for the latter in  %we use the opportunity to deduce the latter from its simpler positive characteristic case via tilting for the \'{e}tale cohomology of perfectoids %strategy used in \cite{brauer-purity} 
%(a perfectoid argument for this result was also noticed by Fujiwara, \cite{Fuj18}). 





%%%%%%%%%%%%%%%%%%%%%%%%%%%%%%%%%%%%%%%%%%%%%%%%%%

\csub[The absolute cohomological purity of Gabber--Thomason] \lab{Gabber-Thomason}


Purity for \'{e}tale cohomology of regular rings, stated precisely in \Cref{abs-coh-pur} (see also footnote~\ref{foot:semipurity}), was conjectured by Grothendieck and settled by Gabber in \cite{Fuj02}, who built on the strategy initiated by Thomason in \cite{Tho84}. Gabber's alternative later proof given in \cite{ILO14}*{expos\'{e} XVI} eliminated the use of algebraic $K$-theory. We present a proof that uses perfectoids, specifically, \Cref{baby-tilting}, to reduce to the positive characteristic case that had been settled by M. Artin already in \cite{SGA4III}*{expos\'{e}~XVI}. The following standard lemmas facilitate the passage to perfectoids.



\blem \label{lem:towers}
\source{purity-flat-cohomology}
Let $(R, \fm)$ be a complete, regular, local ring with residue field $k$.
\benum
\m \label{tower-a}
\source{purity-flat-cohomology}
There is a filtered direct system $\{(R_i, \fm_i)\}_{i \in I}$ of regular, local, finite, flat $R$-algebras that are unramified if so is $R$ \up{see \uS\uref{conv}} such that $\fm_i = \fm R_i$ %each $R_i/\fm_i$ is a finite subextension of $\ov{k}/k$, 
and $(\varinjlim_i R_i, \varinjlim_i \fm_i)$ is a regular local ring whose residue field is an algebraic closure $\ov{k}$ of $k$. 

\m \label{tower-b}
\source{purity-flat-cohomology}
If $R$ is of mixed characteristic $(0, p)$ and $k$ is perfect, then there is a tower $\{R_n\}_{n \ge 0}$ of regular, local, finite, flat $R$-algebras of $p$-power rank over $R$ such that the $p$-adic completion of $R_\infty \ce \varinjlim_{n \ge 0} R_n$ is perfectoid\ucolon explicitly, by the Cohen structure theorem, we have
\[%\be \lab{Rinfty-exp}
\qq R \simeq W(k)\llb x_1, \dotsc, x_d\rrb/(p - f), 
\]
where either $f = x_1$ or $f \in (p, x_1, \dotsc, x_d)^2$ \up{the two cases correspond to whether or not $R$ is unramified}, %in the sense that $p \in \fm \setminus \fm^2$}, \
and one may choose
\[
\qq \tst R_n \ce W(k)\llb x_1^{1/p^n}, \dotsc, x_d^{1/p^n} \rrb/(p - f)
\]
with
\[
\tst \qq R_\infty \simeq W(k)\llb x_1^{1/p^\infty}, \dotsc, x_d^{1/p^\infty} \rrb/(p - f).
\]
\eenum
\elem

For our somewhat nonstandard use of the notation $\llb \cdot \rrb$ in part \ref{tower-b} above, see \eqref{nst-terminology}.


%We recall and stress that we use the (slightly nonstandard) notation
%\[
%A\llb x_1^{1/p^\infty}, \dotsc, x_d^{1/p^\infty} \rrb = \varinjlim_m A\llb x_1^{1/p^m}, \dotsc, x_d^{1/p^m} \rrb.
%\]

\bpf
In essence, the claims are restatements of \cite{brauer-purity}*{Lemmas 5.1 and 5.2}: part \ref{tower-a} follows from \cite{brauer-purity}*{Lemma 5.1} and its proof, whereas part \ref{tower-b} follows from \cite{brauer-purity}*{Lemma 5.2} and its proof. For a prismatic point of view on the construction of $R_\infty$, see \cite{BS19}*{Remark 3.11}.
\epf




\blem \label{lem:et-to-hat}
\source{purity-flat-cohomology}
Let $A$ be a ring, let $a \in A$ be such that $A$ is $a$-Henselian and has bounded $a^\infty$-torsion, %along the principal ideal generated by an $a \in A$ with $A[a^\infty] = A[a^N]$ for some $N > 0$, 
let $\wh{A}$ be the $a$-adic completion of $A$, and let $\Spec(A[\f 1a]) \subset U \subset \Spec(A)$ be an open. We have
\[
R\Gamma_\et(U, \sF) \isomto R\Gamma_\et(U_{\wh{A}}, \sF)
\]
for every torsion abelian sheaf $\sF$ on $U_\et$. In particular, for every closed subset $Z \subset \Spec(A/aA)$ and every torsion abelian sheaf $\sF$ on $A_\et$, we have
\[
R\Gamma_Z(A, \sF) \isomto R\Gamma_Z(\wh{A}, \sF),
\]
and for a Noetherian ring $R$, an ideal $I \subset R$ such that $R$ is  $I$-Henselian,  and the $I$-adic completion~$\wh{R}$,
\[
R\Gamma_I(R, \sF) \isomto R\Gamma_I(\wh{R}, \sF) 
\]
for every torsion abelian sheaf $\sF$ on $R_\et$.
\elem

\bpf
The claims are special cases of \cite{BC19}*{Theorem 2.3.4, Corollary 2.3.5 (e)}, although we could also use earlier references \cite{Fuj95}*{Corollary~6.6.4} or \cite{ILO14}*{expos\'{e} XX, section 4.4}; %, or \cite{BM20}*{Cor.~1.18}; 
see also \Cref{cor:excision} below.
\epf




\bthm \lab{abs-coh-pur}
For a regular local ring $(R, \fm)$ and a commutative, finite, \'{e}tale $R$-group $G$ whose order is invertible in $R$, 
\[
H^i_\fm(R, G) \cong 0 \q \text{for} \q i < 2 \dim(R).
\]
\ethm

\bpf
We use the local-to-global spectral sequence \cite{SGA4II}*{expos\'{e}~V, proposition~6.4} to assume that $R$ is strictly Henselian and then that $G \simeq \bZ/\ell\bZ$ for a prime $\ell$. We then use \Cref{lem:et-to-hat} to assume 
%\Cref{cor:excision} (or already \eqref{eqn:Fujiwara-Gabber}) then allows us to assume 
that $R$ is also complete. Thus, by the Cohen structure theorem \cite{Mat89}*{Theorem~29.7}, if $R$ is equicharacteristic, then $R \simeq k\llb x_1, \dotsc, x_d\rrb$ for a field $k$ and, by \Cref{lem:et-to-hat} again, we may assume that $R$ is the Henselization of $\bA^d_k$ at the origin. For this $R$ the claim was settled already in \cite{SGA4III}*{expos\'{e}~XVI, th\'{e}or\`{e}me~3.7}, so from now on we assume that our complete, strictly Henselian $R$ is of mixed characteristic $(0, p)$.

Since multiplication by $p$ is an automorphism of $\bZ/\ell\bZ$, the trace map \cite{SGA4III}*{expos\'{e}~XVII, sections~6.3.13 et 6.3.14, proposition 6.3.15~(iv)} allows us to replace $R$ by any module-finite, flat $R$-algebra $R'$ of $p$-power rank over $R$ such that $R'$ is a regular local ring. Thus, by \Cref{lem:towers}~\ref{tower-a} and a limit argument, we may pass to a tower to reduce to the case  when the residue field $k$ of $R$ is algebraically closed (we use \Cref{lem:et-to-hat} to complete $R$ again). We then likewise use \Cref{lem:towers}~\ref{tower-b} to  reduce to showing that
\[
H^{i}_{(p,\, x_1,\, \dotsc,\, x_d)}(R_\infty, \bZ/\ell\bZ) \cong 0 \qxq{for}  i < 2d 
\]
with
\[
R_\infty \cong W(k)\llb x_1^{1/p^\infty}, \dotsc, x_d^{1/p^\infty} \rrb/(p - f), \ \  f\in \fm^2 \cup \{ x_1 \},
\]
knowing that the $p$-adic completion $\wh{R}_\infty$ of $R_\infty$ is perfectoid. The tilt $\wh{R}_\infty^\flat$ of $\wh{R}_\infty$ reviewed in \eqref{monoid-id} is the $\ov{f}$-adic completion of $k\llb (x_1^{\flat})^{1/p^\infty}, \dotsc, (x_d^{\flat})^{1/p^\infty} \rrb$ for some $\ov{f} \in k\llb (x_1^{\flat})^{1/p^\infty}, \dotsc, (x_d^{\flat})^{1/p^\infty} \rrb$: explicitly, 
\[
%\xymatrix@C=80pt{ %k\llb x_1^{1/p^\infty}, \dotsc, x_d^{1/p^\infty}\rrb \ce 
\wh{R}_\infty^\flat \cong \varprojlim\limits_{z \mapsto z^p} \p{k\llb x_1^{1/p^\infty}, \dotsc, x_d^{1/p^\infty}\rrb/(f)},
\]
which via the map $(z_n)_{n \ge 0} \mapsto (z_n^{p^n})_{n \ge 0}$ is identified with
\[
\tst\invlim\limits_n \p{k\llb (x_1^{\flat})^{1/p^\infty}, \dotsc, (x_d^{\flat})^{1/p^\infty}\rrb /(\ov{f}^{p^n})},  
\]
where $x_i^\flat$ corresponds to the $p$-power compatible system $(x_i^{1/p^n})_{n \ge 0}$. Thus, by \eqref{eqn:baby-tilting-supports} with \Cref{lem:et-to-hat} (the latter removes the $\ov{f}$-adic completion), we are reduced to showing 
\[
H^{i}_{(x_1^\flat, \dotsc, x_d^\flat)}(k\llb (x_1^{\flat})^{1/p^\infty}, \dotsc, (x_d^{\flat})^{1/p^\infty}\rrb, \bZ/\ell\bZ) \cong 0 \qxq{for} i < 2d.
\]
By the perfection-invariance of \'{e}tale cohomology, we may replace 
\[
k\llb (x_1^{\flat})^{1/p^\infty}, \dotsc, (x_d^{\flat})^{1/p^\infty}\rrb \qxq{by} k\llb x_1^{\flat}, \dotsc, x_d^{\flat}\rrb,
\]
which brings us to the already discussed equal characteristic case.
%: explicitly, by the Cohen structure theorem, we have
%\be \lab{Rinfty-exp}
%R \cong W(k)\llb x_1, \dotsc, x_d\rrb/(p - f), \q \x{where either} \q  f = x_1 \q \x{or} \q f \in (p, x_1, \dotsc, x_d)^2
%\ee
%(the two cases correspond to whether or not $p \in \fm \setminus \fm^2$), and
%\[
%\tst  \ce  \varinjlim_{m > 0} \p{ W(k)\llb x_1^{1/p^m}, \dotsc, x_d^{1/p^m} \rrb/(p - f)}.
%\]
\epf




\brems
\remi
As Teruhisa Koshikawa pointed out to us, the above argument also reduces the full absolute cohomological purity for \'{e}tale cohomology, namely, the statement that for all $n \in \bZ_{> 0}$ invertible in $R$ the \'{e}tale-sheafification $\cH^i_\fm$ of the cohomology with supports $H^i_\fm$ satisfies
\[
\qq \cH^i_\fm(-, \bZ/n\bZ) \cong \begin{cases} 0, & \x{for $i \neq 2\dim(R)$,} \\ \underline{\bZ/n\bZ}(-\dim(R)), & \x{for $i = 2\dim(R)$,} \end{cases}
\]
to positive characteristic. Indeed, the isomorphism in degree $i = 2\dim(R)$ is induced by the cycle class map 
\[
\qqq \underline{\bZ/n\bZ}(-\dim(R)) \ra \cH^{2\dim(R)}_\fm(-, \bZ/n\bZ),
\]
which one first argues to be injective as in \cite{Fuj02}*{Lemma 2.3.1}. The bijectivity then becomes the matter of bounding the nonzero stalk of the target, which may be done after passing to $\wh{R}^\flat_\infty$. The vanishing in degrees $i \neq 2\dim(R)$ reduces to positive characteristic as in the proof of \Cref{abs-coh-pur}.

\remi
Another way to pass to the tilt, without using \Cref{baby-tilting}, is to use diamonds developed in \cite{Sch17}. Namely, we consider the ``punctured adic spectrum'' of $R_\infty$ defined as
\[
\tst \qqq U^\ad_{R_\infty} \ce \Spa(R_\infty, R_\infty) \setminus \{ x_1 = \dotsc = x_d = 0 \}
\]
and isomorphic to
\[
\tst \qqq \bigcup_{i = 1}^{d} \Spa\p{R_\infty\langle\f{x_1,\, \dotsc,\, x_d}{x_i}\rangle, R_\infty\langle\f{x_1,\, \dotsc,\, x_d}{x_i}\rangle^+}\!,
\] 
where $R_\infty$ is endowed with its $(x_1, \dotsc, x_d)$-adic topology, so that $U^\ad_{R_\infty}$ is an analytic adic space over $\bZ_p$ (to simplify we ignore the issue of showing that the appearing Huber pairs are sheafy).\footnote{In fact, by \cite{Nak20}*{Example~2.1.4}, each $\Spa$ appearing in the union on the right side is even an affinoid perfectoid.} Likewise, we endow $\wh{R}_\infty^\flat$ with the $(x_1^\flat, \dotsc, x_d^\flat)$-adic topology and consider 
\[
\tst\qqq U^\ad_{\wh{R}_\infty^\flat} \ce \Spa(\wh{R}_\infty^\flat, \wh{R}_\infty^\flat) \setminus \{ x_1^\flat = \dotsc = x_d^\flat = 0 \},
\] 
which is isomorphic to
\[
\tst \qqq \bigcup_{i = 1}^{d} \Spa\p{\wh{R}_\infty^\flat\langle\f{x_1^\flat,\, \dotsc,\, x_d^\flat}{x_i^\flat}\rangle, \wh{R}_\infty^\flat\langle\f{x_1^\flat,\, \dotsc,\, x_d^\flat}{x_i^\flat}\rangle^+}\!
\]
and is a perfectoid space because the coordinate rings of the appearing affinoids inherit perfectness from $\wh{R}_\infty^\flat$ (see \cite{Sch17}*{Proposition 3.5}). By tilting %establishes an equivalence between perfectoid $\wh{R}_\infty$-algebras $A$ and perfectoid $\wh{R}_\infty^\flat$-algebras $A^\flat$ 
(see \Cref{tilt-summary} and, for compatibility of definitions, \cite{BMS18}*{Lemma~3.20}), the universal property of adic localization and \eqref{monoid-id-pi} show that giving a map from a perfectoid space to 
\[
\tst\qqq \Spa\p{R_\infty\langle\f{x_1,\, \dotsc,\, x_d}{x_i}\rangle, R_\infty\langle\f{x_1,\, \dotsc,\, x_d}{x_i}\rangle^+}
\]
amounts to giving a map from its tilt to the perfectoid space 
\[
\tst\qqq\Spa\p{\wh{R}_\infty^\flat\langle\f{x_1^\flat,\, \dotsc,\, x_d^\flat}{x_i^\flat}\rangle, \wh{R}_\infty^\flat\langle\f{x_1^\flat,\, \dotsc,\, x_d^\flat}{x_i^\flat}\rangle^+},
\]
compatibly with overlaps of such rational subsets. Thus, the perfectoid space $U^\ad_{\wh{R}_\infty^\flat}$ represents the diamond associated  by \cite{Sch17}*{Definition~15.5} to the analytic adic space $U^\ad_{R_\infty}$.  Consequently, since the functor that sends an analytic adic space over $\bZ_p$ to its associated diamond induces an equivalence of \'{e}tale sites \cite{Sch17}*{Lemma~15.6}, we obtain the key
\[
\qqq H^i_\et(U_{R_\infty}^\ad, \bZ/\ell\bZ) \cong H^i_\et(U_{\wh{R}_\infty^\flat}^\ad, \bZ/\ell\bZ). %\cong H^i_\et(U_{R_\infty^\flat}, \bZ/\ell\bZ).
\]
%To pass to the tilt as promised, 
It remains to set 
\[
\qqq U_{R_\infty} \ce \Spec(R_\infty) \setminus \{ x_1 = \dotsc = x_d = 0\}
\]
and
\[
\qqq U_{R^\flat_\infty} \ce \Spec(R^\flat_\infty) \setminus \{ x^\flat_1 = \dotsc = x^\flat_d = 0\}
\]
and apply \cite{Hub96}*{Theorem~3.2.10}\footnote{Due to the blanket Noetherianity assumption  \cite{Hub96}*{paragraph~1.1.1}, the citation does not apply directly, so one performs a slightly tedious limit argument, similar to the one used in the proof of \cite{brauer-purity}*{Theorem 4.10}. Another way around this is to pass to the adic spectra at the finite levels of the tower and then use \cite{Sch17}*{Proposition 14.9} to pass to the limit of adic spaces.} (for the flanking isomorphisms) to deduce the passage to the tilt:
\[
\qq\ \ H^i_\et(U_{R_\infty}, \bZ/\ell\bZ) \cong  H^i_\et(U_{R_\infty}^\ad, \bZ/\ell\bZ) \cong H^i_\et(U_{\wh{R}_\infty^\flat}^\ad, \bZ/\ell\bZ) \cong H^i_\et(U_{R_\infty^\flat}, \bZ/\ell\bZ).
\]
%where , and the first and last isomorphisms are supplied by  %(comparison between the algebraic and the analytic \'{e}tale cohomology).
\erems





%\blem \lab{lem:et-GAGA}
%For a ring $A$ that is Henselian along a finitely generated ideal $I \subset A$ and is endowed with the $I$-adic topology, the opens $U \ce \Spec(A) \setminus V(I)$ and $U^\ad \ce \Spa(A, A) \setminus V(I)$ satisfy\footnote{\rv{\me{There is some hairy issue of how the \'{e}tale site $U^\ad_\et$ is defined. Apparently, one should use the same definition as in \cite{KL15}*{8.2.19}.}}}
%\[
%R\Gamma_\et(U, G) \isomto R\Gamma_\et(U^\ad, G) \q \x{for any commutative finite \'{e}tale $U$-group scheme $G$.}
%\]
%\elem

%\emph{Proof.} 
%In the case when $A$ is Noetherian, the claim is a special case of . In general, we \kestutis{unfinished} \QEDD
















%%%%%%%%%%%%%%%%%%%%%%%%%%%%%%%%%%%%%%%%%%%%%%%%%%

\csub[The \'{e}tale depth is at least the virtual dimension] \lab{et-depth-geom-depth}



Purity for \'{e}tale cohomology of possibly singular Noetherian local rings $R$ was settled in the case when $R$ is an excellent $\bQ$-algebra in \cite{SGA2new}*{expos\'{e} XIV, th\'{e}or\`{e}me 5.6} and in the case when $R$ is a complete intersection in \cite{Gab04}*{Theorem~3} (by reduction to \cite{Ill03}*{th\'{e}or\`{e}me~2.6}). In \Cref{ell-semipurity}, we deduce the general case from \Cref{abs-coh-pur}. We begin with the definition of the virtual dimension, which is a numerical invariant of $R$ that has already appeared in the context of purity in \cite{SGA2new}*{expos\'{e} XIV, d\'{e}finition 5.3}.
%The \'{e}tale case of our desired purity says that the ``\'{e}tale depth'' of a Noetherian local ring $(R, \fm)$, concretely, the smallest $d$ such that the \'{e}tale cohomology $H^i_\fm(R, G)$ with commutative, finite, \'{e}tale $R$-group coefficients $G$ vanishes for $i < d$, is at least the virtual dimension $\vdim(R)$. For $G$ whose order is invertible in $R$, we deduce this from  in \Cref{ell-semipurity} below. 









\bpp[The virtual dimension of a Noetherian local ring] \lab{geom-depth-def}
For a Noetherian local ring $(R, \fm)$, by Cohen's theorem \cite{EGAIV1}*{chapitre 0, th\'{e}or\`{e}me 19.8.8 (i)}, the $\fm$-adic completion $\wh{R}$ is of the form
\[%be \lab{Cohen-pres}
\wh{R} \cong \wt{R}/I \qq \text{for a complete regular local ring} \q \wt{R} \q \text{and an ideal} \q I \subset \wt{R}.
\]%ee
The \emph{virtual dimension} of $R$ is %defined by
\be \label{eqn:vdim-def}
\source{purity-flat-cohomology}
\vdim(R) \ce \dim(\wt{R}) - (\text{minimal number of generators for $I$})
\ee
and, by \cite{SGA2new}*{expos\'{e} XIV, proposition 5.2}, does not depend on the presentation $\wt{R}/I$. %\eqref{Cohen-pres} (in \emph{op.~cit.}~this invariant is called the \emph{geometric depth} of $R$). 
By \cite{SGA2new}*{expos\'{e} XIV, proposition 5.4}, %we have
%Following \cite{SGA2new}*{XIV, 5.3}, we define the \emph{virtual dimension} of $R$ by
\be \label{vdim-ineq}
\source{purity-flat-cohomology}
\vdim(R) \le \dim(R) \q \text{with equality iff $R$ is a complete intersection.}
\ee
By definition, $\vdim(R) = \vdim(\wh{R})$, so also $\vdim(R) = \vdim(R^h)$; %in addition, $\vdim(R) = \vdim(R^\sh)$, in fact,
%the virtual dimension is invariant under replacing $R$ by its completion, so also by its Henselization. The same holds for the strict Henselization: in fact, 
more generally, by \cite{Avr77}*{Proposition 3.6 and equations (3.2.1) and (3.2.2)}, for any flat local homomorphism $R \ra R'$ of Noetherian local rings, we have
\be \label{vdim-add}
\source{purity-flat-cohomology}
\vdim(R') = \vdim(R) + \vdim(R'/\fm R'), 
\ee
so, in particular,
\[
\vdim(R) = \vdim(R^{\sh})
\]
(\emph{loc.~cit.}~proves this for the \emph{complete intersection defect} defined as $\dim(*) - \vdim(*)$ but the dimension $\dim(*)$ is likewise additive, see \cite{EGAIV2}*{corollaire~6.1.2}). %, which extends \cite{SGA2new}*{XIV,~5.5}.
\epp


\brem
Despite the name ``geometric depth'' used for $\vdim(R)$ in \cite{SGA2new}*{expos\'{e}~XIV, d\'{e}finition 5.3}, in general there is no inequality between $\depth_\fm(R)$ and $\vdim(R)$: a Cohen--Macaulay $R$ that is not a complete intersection has $\depth_\fm(R) > \vdim(R) $, whereas, due to \cite{Bur68} (or \cite{Koh72}*{Theorem~A}) and the Auslander--Buchsbaum formula, any regular local ring $\wt{R}$ has an ideal $I$ generated by three elements such that the quotient 
\[
R \ce \wt{R}/I \qxq{has} \depth_\fm(R) = 0
\]
(and $\vdim(R) \ge \dim(\wt{R}) - 3$). % MO #278191
\erem




To deduce \Cref{ell-semipurity} from \Cref{abs-coh-pur}, we will use the following Lefschetz hyperplane theorem in local \'{e}tale cohomology. This strategy is close in spirit to the one used by Mich\`{e}le Raynaud in \cite{SGA2new}*{expos\'{e} XIV, th\'{e}or\`{e}me 5.6} in the case when $R$ is an excellent $\bQ$-algebra. %Other methods are possible: for instance, since we first reduce to complete intersection $R$, we could afterwards use a reduction indicated in  to arrive at the case when $R$ is essentially of finite type over a discrete valuation ring that had been settled in .




%The prime to the characteristic aspects of our main result, \Cref{main}, are not new even though their proofs in the stated generality do not seem to appear in the literature. For instance, their complete intersection case was announced in 

% The argument we give here (in the proof of \Cref{ell-semipurity}) is closer to the one for , which is the case of an excellent $\bQ$-algebra: we take advantage of numerous improvements due to Gabber \cite{ILO14} to the more restrictive earlier results from \cite{SGA4III} and \cite{SGA5} in order to eventually deduce the conclusion from the absolute cohomological purity theorem. The statement of \Cref{ell-semipurity} involves the virtual dimension, so we begin with a 


\blem \label{lem:local-hyperplane}
\source{purity-flat-cohomology}
For a regular local ring $(R, \fm)$, an $f \in \fm$, and an invertible in $R$ prime $\ell$, the map
\[
\q H^i_\fm(R, \bZ/\ell\bZ) \ra H^i_\fm(R/(f), \bZ/\ell\bZ) \qxq{is} \begin{cases} \x{bijective for} \q i < \dim(R) - 1, \\ \x{injective for} \q i = \dim(R) - 1. \end{cases}
\]
\elem

\bpf
Letting $j \colon \Spec(R[\f{1}{f}]) \hra \Spec(R)$ be the indicated open immersion, we need to show that
\be \label{eqn:j!-vanishes}
\source{purity-flat-cohomology}
H^i_{\fm}(R, j_!(\bZ/\ell \bZ)) \cong 0 \q \x{for} \q i < \dim(R).
\ee
Moreover, by the local-to-global spectral sequence \cite{SGA4II}*{expos\'{e}~V, proposition~6.4}, we may assume that $R$ is strictly Henselian and, by \Cref{lem:et-to-hat}, that $R$ is also complete. We will derive \eqref{eqn:j!-vanishes} from Gabber's affine Lefschetz theorem \cite{ILO14}*{expos\'{e} XV, corollaire 1.2.4}, which gives
\be \lab{aff-Lef}
\tst H^i(R[\f{1}{f}], \bZ/\ell\bZ) \cong 0 \q \x{for} \q i > \dim(R).
\ee
Namely, by %Gabber's results on dualizing complexes 
\cite{ILO14}*{expos\'{e} XVII, th\'{e}or\`{e}me 0.2}, the complex of \'{e}tale sheaves $(\mu_\ell^{\tensor \dim(R)})[2\dim(R)]$  on $R$ is dualizing and its $j^!$-pullback is dualizing on $R[\f{1}{f}]$. This pullback is $(\mu_\ell^{\tensor \dim(R)})[2\dim(R)]$ (see \cite{SGA4III}*{expos\'{e}~XVIII, proposition~3.1.8~(iii)}), so, since 
\[
\bR j_* \circ \bD \cong \bD \circ j_!
\]
(see \cite{SGA5}*{expos\'{e}~I, proposition~1.12~(a)}), the vanishing \eqref{aff-Lef} amounts~to
\[
\tst \p{H^{i}\p{\bR\sH om\p{j_!(\bZ/\ell \bZ), (\mu_\ell^{\tensor \dim(R)})[2\dim(R)]}}}_{\fm} \cong 0 \q \text{for} \q i > -\dim(R),
\]
where $(-)_{\fm}$ indicates the stalk. To obtain \eqref{eqn:j!-vanishes} it remains to use local duality in \'{e}tale cohomology \cite{SGA5}*{expos\'{e}~I, \'{e}quation~(4.2.2)} (our dualizing complex is normalized %at $\fm$ 
as there by \cite{ILO14}*{expos\'{e}~XVI, th\'{e}or\`{e}me~3.1.1}). %or \cite{Fuj02}*{2.1.1}). %the vanishing \eqref{inter-van} amounts to the desired .
\epf















%\temp{
%The following simple lemma captures a crucial trick that goes into the proofs of Theorems \ref{main} and \ref{ell-semipurity}: roughly, a well-chosen $\varpi$ can play the role of $p$ in the perfectoid parts of the overall argument---this bypasses all complications caused by $p$-torsion in mixed characteristic perfectoids.

%\blem \label{lem:choose-varpi}
%A Cohen--Macaulay local ring $(R, \fm)$ of residue characteristic $p > 0$ and positive dimension has a regular sequence $\varpi, r_2, \dotsc, r_{\dim(R)} \in \fm$ such that $R/(\varpi)$ is killed by a power of $p$. 
%\elem

%\bpf
%Let $N > 0$ be large such that $p^N$ vanishes at the localizations of $R$ at the minimal primes of residue characteristic $p$. 
%By prime avoidance \cite{SP}*{\href{https://stacks.math.columbia.edu/tag/00DS}{00DS}}, there is an $f \in \fm$ that does not vanish at the minimal primes of residue characteristic $p$ but vanishes at every minimal prime of residue characteristic $0$. We replace $f$ by a large power to ensure that $f$ even maps to $0$ in $R_\fp$ for every minimal prime $\fp \subset R$ of residue characteristic $0$. Then, since Cohen--Macaulay rings have no embedded associated primes, $\varpi \ce p + f$ is a nonzerodivisor whose image in $R_\fq$ for every prime $\fq$ of residue characteristic $0$ is $p$. Consequently, $R/(\varpi)$ is Cohen--Macaulay and all of its primes are of residue characteristic $p$. Thus, $R/(\varpi)$ is killed by a power of $p$, and $\varpi$ may be extended to a desired regular sequence.
%\epf
%}





\bthm \lab{ell-semipurity}
For a Noetherian local ring $(R, \fm)$ and a commutative, finite, \'{e}tale $R$-group $G$ whose order is invertible in $R$,
\[
H^i_\fm(R, G) \cong 0 \q \text{for} \q i < \vdim(R).
\]
\ethm

%We will remove the unnecessary assumption on the order of $G$ in the proof of \Cref{main-etale} (see~\S\ref{main-etale-pf}).

We will remove the assumption on the order of $G$ in \Cref{ell-semipurity-p} below.

\bpf
\Cref{abs-coh-pur} and \Cref{lem:local-hyperplane} settle the case when $R$ is a hypersurface, that is, a quotient of a regular ring by a principal ideal. Thus, it suffices to show how to reduce from a general $R$ to a hypersurface. This reduction works for any commutative, finite, \'{e}tale $G$, so, to be able to reuse it in the proof of \Cref{ell-semipurity-p}, we drop the assumption on the order of $G$.

By the spectral sequence \cite{SGA4II}*{expos\'{e}~V, proposition~6.4}, we may assume that $R$ is strictly Henselian, and then, by d\'{e}vissage, that $G \simeq \bZ/p\bZ$ for a prime $p$. \Cref{lem:et-to-hat} reduces further to complete $R$, so that %
\[
R \simeq \wt{R}/(f_1, \dotsc, f_n)
\]
for a complete regular local ring $(\wt{R}, \wt{\fm})$ and $f_1, \dotsc, f_n \in \wt{\fm}$ chosen so that $f_1, \dotsc, f_n$ is a minimal generating set for the ideal $(f_1, \dotsc, f_n) \subset \wt{R}$ (see \S\ref{conv}). Suppose that $n > 1$ and consider the rings 
\[
R_1 \ce \wt{R}/(f_1, \dotsc, f_{n - 1}) \qxq{and} R_2 \ce \wt{R}/(f_n),
\]
as well as 
\[
R_{12} \ce \wt{R}/(f_1f_n, \dotsc, f_{n - 1}f_n).
\]
Set theoretically, in $\Spec(\wt{R})$ we have 
\[
\Spec(R_1) \cap \Spec(R_2) = \Spec(R) \q \x{and} \q \Spec(R_1) \cup \Spec(R_2) = \Spec(R_{12}).
\]
Thus, since the \'{e}tale site is insensitive to nilpotents, we obtain the exact sequence
\[
0 \ra \underline{\bZ/p\bZ}_{R_{12}} \xra{z\, \mapsto\, (z|_{R_1},\, z|_{R_2})} \underline{\bZ/p\bZ}_{R_1} \oplus \underline{\bZ/p\bZ}_{R_2} \xra{(x,\, y)\, \mapsto\, x|_R - y|_{R}} \underline{\bZ/p\bZ}_{R} \ra 0
\]
of sheaves on $\Spec(\wt{R})_\et$. Since $\vdim(R_*) \ge \vdim(R) + 1$ for $* \in \{1, 2, 12\}$, the associated long exact sequence of cohomology with supports reduces the desired vanishing to its counterpart for the rings $R_*$. This allows us to decrease $n$, so  we arrive at $n = 1$, that is, at $R$ being a hypersurface $\wt{R}/(f)$. %so that $\vdim(R) = \dim(R)$. % applied to $\wt{R}$ then give the case when $p$ is invertible in $R$. 
\epf


%Other ways to conclude: via tilting and \cite{Ill05}*{Thm.~1.3} or via a limit argument and \cite{Ill05}*{Thm.~2.6}.









%%%%%%%%%%%%%%%%%%%%%%%%%%%%%%%%%%%








\csub[Zariski--Nagata purity for rings of virtual dimension $\ge 3$] \lab{nonabelian}



The nonabelian analogue of \Cref{ell-semipurity} is the following generalization of the Grothendieck--Zariski--Nagata purity theorem \cite{SGA2new}*{expos\'{e} X, th\'{e}or\`{e}me 3.4} (and of the main result of \cite{Cut95}): for extending finite \'{e}tale covers over a  closed subscheme of a Noetherian scheme, it suffices to assume that the total space have virtual dimension $\ge 3$ at the missing points (instead of even being a complete intersection of dimension $\ge 3$ at these points).
 %the complete intersections of dimension $\ge 3$ the complete intersection assumption is not needed,  suffices. 
 We learned from de Jong that his generalizations contained in \cite{SP} of the algebraization theorems from \cite{SGA2new} and \cite{Ray75} could be used to prove this---indeed, as the reader will notice, they are the main inputs to the proof of \Cref{non-ab-et}. This section is purely classical and does not use perfectoid inputs or other sections. 


\bthm \lab{non-ab-et}
Let $(R, \fm)$ be a Noetherian local ring, set $U_R \ce \Spec(R) \setminus \{ \fm\}$, and let $\Spec(R)_\fet$ \up{resp.,~$(U_{R})_\fet$} denote the category of finite \'{e}tale $R$-schemes \up{resp.,~of finite \'{e}tale $U_{R}$-schemes}.
\benum
\item \lab{NAE-a}
If $\vdim(R) \ge 2$, then the pullback $\Spec(R)_\fet \ra (U_{R})_\fet$ is fully faithful, $U_R$ is connected, and 
\[
\q \pi_1^\et(U_R) \surjects \pi_1^\et(R).
\]

\item \lab{NAE-b}
If $\vdim(R) \ge 3$, then the pullback $\Spec(R)_\fet \ra (U_{R})_\fet$ is an equivalence of categories and
\[
\q \pi_1^\et(U_R) \isomto \pi_1^\et(R).
\]
\eenum
\ethm


\bpf
In \ref{NAE-a}, granted the full faithfulness, the connectedness of $U_R$ follows from that of $\Spec(R)$  by considering sections both of the finite \'{e}tale map $\Spec(R) \bigsqcup \Spec(R) \ra \Spec(R)$ and of its base change to $U_R$. Moreover, by \cite{SGA1new}*{V, 6.9, 6.10}, % another reference is \cite{Sza09}*{5.5.4 2. and 5.5.8}, although it does not seem as convenient (no direct relation to full faithfulness)
%in both \ref{NAE-a} and \ref{NAE-b} 
the conclusions about the fundamental groups follow from the claims about the functors. For the latter, patching \cite{FR70}*{proposition~4.2} and flat descent allow us to replace $R$ by its $\fm$-adic completion. Then we may write 
\[
R \simeq \wt{R}/(f_1, \dotsc, f_n)
\]
for a complete regular local ring $(\wt{R}, \wt{\fm})$ and $f_1, \dotsc, f_n \in \wt{\fm}$ chosen so that $f_1, \dotsc, f_n$ is a minimal generating set for the ideal $\ff \ce (f_1, \dotsc, f_n) \subset \wt{R}$ (see \S\ref{conv}). Let $\fU \subset \fX$ be the formal schemes obtained from 
\[
U_{\wt{R}} \ce \Spec(\wt{R}) \setminus \{ \wt{\fm}\} \subset  \Spec(\wt{R})
\]
 by formal $\ff$-adic completion. %, and consider the Zariski open $\fU \subset \fX$ obtained by formally $\ff$-adically completing . 
Since  \'{e}tale sites are insensitive to nilpotents, pullback gives equivalences of categories
\[
\fX_\fet \isomto \Spec(R)_\fet \qq \x{and} \qq \fU_\fet \isomto (U_{R})_\fet,
\]
so we are reduced to considering the pullback functor $\fX_\fet \ra \fU_\fet$. By \cite{SP}*{Lemma~\href{https://stacks.math.columbia.edu/tag/09ZL}{09ZL}}, %the Grothendieck algebraization theorem \cite{EGAIII1}*{5.1.4}, %to get flatness: finite \'{e}tale maps to a local ring are free as modules and the trivializations may be lifted due to Nakayama, so we are really only algebraizing free modules. 
the pullback $\Spec(\wt{R})_\fet \isomto \fX_\fet$ is an equivalence, so we only need to consider the pullback $\Spec(\wt{R})_\fet \ra \fU_\fet$.

\benum
\item
The assumption $\vdim(R) \ge 2$ allows us to apply the algebraization of formal sections \cite{SP}*{Lemma~\href{https://stacks.math.columbia.edu/tag/0DXR}{0DXR} (with Lemma~\href{https://stacks.math.columbia.edu/tag/0DX9}{0DX9})} %one applies this to Hom sheaves, which are locally free, so their formation commutes with pullback
(or \cite{Ray75}*{chapitre~I, th\'{e}or\`{e}me~3.9}) to conclude that any morphism between the $\fU$-pullbacks of finite \'{e}tale $\wt{R}$-schemes $Y_1$ and $Y_2$ algebraizes to a morphism between their $U$-pullbacks for some open $U_R \subset U \subset U_{\wt{R}}$, and this algebraization is unique up to shrinking $U$. The complement 
\[
\qq Z \ce \Spec(\wt{R}) \setminus U
\] 
is at most $n$-dimensional because the $f_i|_Z$ cut out the closed point of $Z$. Thus, since $n \le \dim(\wt{R}) - 2$, the codimension of $Z$ in $\Spec(\wt{R})$ is $\ge 2$, to the effect that the algebraized morphism extends uniquely to an $\wt{R}$-morphism between $Y_1$ and $Y_2$ (see \cite{EGAIV2}*{th\'{e}or\`{e}me 5.10.5}). Consequently, the pullback $\Spec(\wt{R})_\fet \ra \fU_\fet$ is fully faithful.

\item
The assumption $\vdim(Rd) \ge 3$ allows us to apply the algebraization of coherent formal modules \cite{SP}*{Lemma~\href{https://stacks.math.columbia.edu/tag/0EJP}{0EJP}} to conclude that any finite \'{e}tale $\fU$-scheme algebraizes to a finite \'{e}tale $U$-scheme $Y$ for some open $U_R \subset U \subset U_{\wt{R}}$ (to algebraize the algebra structure maps, we use \cite{SP}*{Lemma~\href{https://stacks.math.columbia.edu/tag/0DXR}{0DXR}} as in \ref{NAE-a}). 
The complement of $U$ in $\Spec(\wt{R})$ is now of codimension $\ge 3$, %(see the proof of \ref{NAE-a}), 
so, by the Zariski--Nagata purity for regular rings \cite{SGA2new}*{expos\'{e} X, th\'{e}or\`{e}me 3.4 (i)}, we may extend $Y$ to a finite \'{e}tale $\wt{R}$-scheme. Consequently, the pullback $\Spec(\wt{R})_\fet \ra \fU_\fet$ is essentially surjective.
\qedhere
\eenum
\epf

\brem
The connectedness of $U_R$ holds more generally, see, for instance, \cite{SP}*{Lemma~\href{https://stacks.math.columbia.edu/tag/0ECR}{0ECR}} or \cite{Var09}*{Theorem~1.6}.
\erem







%An application in the spirit of \cite{Cut95}*{Thm.~22} should be possible: if $R$ is (R$_2$) and of virtual dimension $\ge 3$, then $\Cl(R)$ has no torsion of order invertible on $R$ (a generalization of \cite{Gri01}*{1.8}). Global version on representing elements of $\Cl(X)$ by line bundles for (R$_2$) $X$ whose local rings of dim $\ge 3$ are of virtual dimension $\ge 3$?





















%%%%%%%%%%%%%%%%%%%%%%%%%%%%%%%%%%%%%%%%%%%%%%
